\section{Discussion}
\label{sec:discussion}

\para{What the evidence supports.}
We tested four hypotheses.
(H1) The memory-based observers beat \directao only on the caution event, by $+0.03$ to $+0.11$ AUROC, with CIs that touch zero; the unsafe event cannot discriminate between methods.
(H2) The gain does not grow with lead distance; it peaks at $H=256$--$512$.
(H3) Among components, memory helps at long horizons ($-0.09$ to $-0.11$ AUROC at $H \ge 512$ when removed from \roll), while the JEPA loss, the hazard head and latent rollout bring no measurable gain.
(H4) Content forecasting beyond trivial baselines is supported: behavior is forecast at 0.87--0.88 accuracy 512--1024 tokens ahead, and summary retrieval reaches about 5$\times$ chance MRR.

\para{Much of what is forecastable is prompt-level.}
Strong simple baselines are a recurring theme.
A linear probe on the latest activation is competitive on the caution event, and text-only TF-IDF is the best unsafe-event predictor.
Much of the behavior signal is response genre (\eg narrate vs.\ explain), which persists through a response.
These observations suggest that, at this scale, the forecastable content is mostly a prompt-level or genre property, readable from the prompt-end activation or prompt text, rather than a long-range dynamic plan.
This is consistent with prior findings that prompt states encode global response properties~\citep{dong2025emergent} and that text-only baselines can rival activation readers~\citep{bauer2026building}.

\para{Why the hazard head did not help.}
In principle a hazard head gives monotone, coherent risk across horizons.
In practice, with 86\% of responses censored by \eos, late bins receive almost no supervision, and the product in \eqref{eq:hazard} extrapolates poorly; its Brier score at $H=2048$ is ten times that of independent classifiers.
A parametric tail, or a mixture of direct per-$H$ heads with a monotonicity constraint, may keep the coherence without the extrapolation error.

\para{Limitations.}
\emph{Scale.} This is a pilot with 84 trajectories, 26 test trajectories in 12 test families, 5 unsafe and 7 late caution events; most CIs span $\pm$0.05--0.2.
\emph{Natural events are early.} No unsafe onset occurs after 448 tokens, so the 1024--2048 horizons are untestable for that event, and long-$H$ labels come from a biased subset of long \longwriter outputs.
\emph{Caution is partly a refusal signal.} Most caution onsets are at token 0; those trajectories contribute only cutoff 0, and late onsets drive the long-$H$ results.
\emph{Labels.} A single LLM annotator gives single-pass labels at 64-token resolution, with no inter-annotator check and no cross-validation against a guard model.
\emph{Baselines.} We did not run an off-the-shelf Activation Oracle or LatentQA decoder; \directao is a matched MLP readout, and our ``decoders'' are classification and retrieval heads rather than free-text generators.
\emph{Other.} We used one rollout per prompt, so we cannot measure calibration against the model's true outcome distribution; hyperparameters were not tuned; and $p$-values are not corrected for multiple comparisons, which would only strengthen the null conclusions.

\para{Safety considerations and broader impact.}
An early design that wrapped harmful requests in long-form templates to construct late unsafe onsets was abandoned, because it would have elicited detailed harmful content.
We use only naturally occurring events from unmodified benchmark prompts, and we release no harmful generations.
Better forecasting monitors could help stop harmful generations early; the same monitors could in principle be used to probe for prompts that delay detection, which is one reason we report their failure modes in detail.
